\chapter{Pioneers}

Computing at UT in the seventies reflected the spirit of the times. This was the era of the Bay Area Homebrew Computer Club at the Stanford Linear Accelerator Center, and the Stanford and MIT AI labs were near their zenith. There was a national undercurrent where government facilities and funding indirectly fostered the growth of hacker culture and community computing, and a certain institutional porosity, where the federal government’s investment in infrastructure like the Southwest Network and Arpanet provided the environment for both rigorous research and independent exploration. When coder Ray Tomlinson introduced electronic mail to the Arpanet in 1971, users quickly developed mailing-list software that allowed a single message to be broadcast to a large group. Almost immediately, this technology was used for socialization. The very first mailing list on the Arpanet was Sf-lovers, created specifically to connect science fiction fans across the country. 

Initially, Arpa administrators frowned upon the network being used for unofficial purposes and explicitly demanded that users dissolve all unauthorized mailing lists. However, the users cleverly pushed back, convincing Arpa to allow the lists by arguing that sending massive amounts of text to Sf-lovers was a necessary way to test the network's mail capacity. Before the Arpanet, brilliant programmers and early hackers were largely confined to isolated labs. The Arpanet's electronic highways effectively merged these localized groups into a single, continent-wide networked tribe. Suddenly, these programmers could easily share information, leading to the propagation of the first shared slang and jargon files, satires, and widespread discussions that codified the modern hacker ethic.

Eventually, Arpa adopted a policy of conscious leniency regarding the network's use for these social lists and early network games. The agency realized that the minimal computational overhead required to host these activities was a highly acceptable cost. By allowing this unofficial playground to exist, the government effectively subsidized the attraction, training, and retention of a generation of creative young computer researchers. This institutional tolerance for playful cleverness is exactly what allowed students at nodes like UT Austin to commandeer university mainframes such as the DEC-10 and transform them from rigid academic tools into shared virtual spaces.

A curious example of all this is the story of Dave Ahl, his computer magazine Creative Computing, and its relationship with events at UT. Creative Computing and Byte magazines were the twin pillars of the industry and friendly competitors. They served different niches and even collaborated occasionally before being absorbed by larger corporate entities. Creative Computing, founded by Ahl in 1974, is widely considered the first personal computer magazine and focused on the educational and playful side of computing. Byte was launched about a year later and quickly became the journal of record for the industry, known for its technical depth and massive, brick-like monthly issues.

Before founding Creative Computing, Ahl was an educational product manager at DEC, and in 1973 he was directly responsible for popularizing the Star Trek game. The first version was written in 1971 by a high school senior named Mike Mayfield. He originally wanted to create a game similar to the graphical space combat game Spacewar, but because he only had access to a non-graphical teletype connected to an SDS Sigma 7 mainframe at UCLA, he was forced to invent a text-based, grid-coordinate approach written in Basic. This UCLA Sigma 7 also happened to be the very first computer on the Arpanet. The first IMP was installed next to it on August 30, 1969, and within a few months it was linked with the second IMP at the Stanford Research Institute.

In 1972, Mayfield rewrote the Star Trek game from scratch for an HP 2000C minicomputer, again in Basic, and submitted it to Hewlett-Packard's public domain library as Sttr1, attributing it to his fictional company Centerline Engineering. This is the version that Ahl and Mary Cole ported to DEC Basic in the summer of 1973. Ahl published the source code in the DEC education-market newsletter, and then under the name Spacwr in his highly influential 1973 book 101 Basic Computer Games, which eventually became the first million selling computer book. This publication is what truly seeded the game across global academic and research institutions, inspiring countless local variations.

Meanwhile, a separate and quirky evolutionary branch was developing at UT. In 1973, Grady Hicks and Jim Korp developed a Basic variant on the CDC 6600. They explicitly acknowledged in their source code that they were building on work already done at the University of Pennsylvania, but apparently did not mention the Ahl or Mayfield branches. This new CDC 6600 version became infamous for a highly disruptive feature. It would periodically halt gameplay to slowly print out lengthy philosophical quotes from Marcus Aurelius's Meditations. Because players were using mechanical teletypes operating at a painfully slow 110 baud, about 10 characters per second, a single quote could consume several minutes of precious machine time and ruin the pacing of the game. The Basic program also took a long time to load and its demands on the busy CDC 6600 irritated the Computation Center staff.

Frustrated by the Marcus Aurelius quotes and the slow loading time, UT grad students Dave Matuszek and Paul Reynolds rewrote the game in 1973 and 1974 using Fortran. Their version loaded quickly and introduced significant gameplay depth, including stars that could go nova and a mysterious enemy called The Thing. This Fortran version eventually evolved into a two-player game called War around 1976, still on the CDC 6600. The two-player version was played on a single teletype, with the players taking turns at the keyboard. Video terminals were becoming more common by 1976 and the fact that a teletype was used indicates how the hardware and environment associated with the CDC 6600 was already becoming somewhat old-fashioned and even obsolete. Since the DEC-10 had already arrived in 1975, video terminals were presumably being used to connect to it, while the old teletypes began to gather dust and were free for two-player Star Trek.

Dave Ahl and Creative Computing  magazine reentered the story at that point with a two-page 1976 article titled "Computing At The University of Texas". UT had recently acquired the DEC-10 specifically because it prioritized interactive time-sharing and research, educational, and in some sense creative uses. This article included a black and white picture of the DEC-10 in its HRC machine room with staff members Pat Caroom, Judith West, and Elissa Vogel. This magazine article would become an important piece of UT history, as it is now easily available on the internet along with most of the Creative Computing content. Creation of the article likely began in 1975, soon after the DEC-10 was installed. With his contacts back to the educational program at DEC, Ahl was likely aware of the new DEC-10 and requested the article, if he didn’t in fact write it himself. The actual authorship is not clear, but with the level of detail in the article it seems likely that it was written at UT and contributed to the magazine. At that point, in 1976, both the original UT Fortran Star Trek game and its Fortran two-player descendant War were still running on the CDC 6600. They were essentially direct descendants or relatives of the original 1971 Star Trek game that Ahl had championed since 1973. And Ahl was now unwittingly publishing an article about the DEC-10 on which, two years later, in 1978, an ultimate version would be created.

The DEC-10 provided a style of timesharing environment that was new to the broader UT community. Operating under the Tops-10 timesharing operating system, it offered a more expansive user experience. With a large 36-bit central memory, still consisting of bulky magnetic cores, and a wide range of interactive peripherals, the DEC-10 allowed users to more easily write, edit, and debug programs in real-time from remote terminals. New levels of interactivity were a catalyst for new types of workflows and research. Meanwhile, it was perfectly true that the 36-bit architecture was well-suited for scientific computing, with 36-bits providing roughly eight decimal digits of precision. And double-precision hardware instructions handled 62-bit fractions for approximately eighteen decimal digits of precision, a feature that was essential for the complex astronomical and engineering models developed at UT. 

An equally important change was taking place at the same time with the proliferation of serial networking, dialup modem access, and the emergence of video terminals in the key terminal rooms, such as the ground-floor facility in Pearce Hall, the old law building, where graduate students like Dave Matuszek worked and provided adhoc technical support to the general users. By 1974, Matuszek had his own video terminal in his office in Pearce Hall, with green text and 80 by 24 screen. He has specifically mentioned using punchcards when he first arrived at UT a few years earlier, and then teletypes, but by 1974 the phosphorescent glow of video terminals had arrived, at least for computer science graduate students. 

\section{Teletypes And Cathode Ray Tubes}

Around 1970, room 105 in Pearce Hall, the old law school, was a vital location for student engagement with the CDC 6600, serving as a primary intake point where punchcard decks were submitted into an input basket for handling by Computation Center staff and an eventual run on the machine. Real-time interactive computing was treated as a scarce privilege. Teletypes were sequestered in locked rooms within the Hogg building, forcing students to navigate administrative friction just to check out a physical key. In the fall of 1971, the computer science department was housed in Waggener Hall together with classics and philosophy. Pearce Hall was to the south, across inner campus drive, just west of the business building. Room 105 was a remote job entry terminal, and where Dave Matuszek remembers going to submit card decks and retrieve printouts. These functions could also be performed in the Computation Center proper.

A year later, in 1972, computer science moved to Painter Hall, north of the Main Building, where it stayed for many years. Painter was shared with the physics department, and also housed the Painter Hall telescope. If any building can be dubbed the home of Star Trek, it would be Painter, just as the HRC was the home of Decwar. In the early eighties there was talk of constructing a new dedicated building for computer science. A slot was even reserved in the construction schedule for this, but much to the dismay of many computer science faculty, this slot was stolen by the new MCC building at Braker and Mopac. The Taylor Hall annex on the corner of 24th and Speedway, across the street from aerospace in WRW, housed another of the Computation Center remote job entry sites, together with several user services and consulting offices. It was torn down in the early nineties to make room for the applied computational engineering and sciences building. By then, the computer science department had moved from Painter into Taylor Hall proper.

The story of Pearce Hall is an example of temporary institutional repurposing. Formerly the law building and later occupied by the anthropology department, it was utilized for computing support precisely because it was obsolete. This allowed university administrators to modify and reconfigure rooms to accommodate the heat and space requirements of a terminal room, without the bureaucratic resistance that would typically accompany such renovations in a permanent facility. This era of makeshift innovation culminated in a terminal room with a jumbled mix of both teletypes and video terminals, confirmed by Dave Matuszek. This was immediately before Pearce Hall’s removal, around 1975, to make way for construction of the new graduate school of business building. Probably the terminal room was one of the last useful jobs for the old building. A very similar story was played out later with the Taylor Hall annex, experimental science building, and WRW.

While the spread of video terminals during the seventies was being catalyzed by companies like Datapoint in San Antonio, their availability on campus was localized and uneven, with the full strength of the video terminal era only truly taking hold around 1974, just as Pearce Hall was being phased out. The transition was guided by the university-wide Faculty Computer Committee. This influential body, which featured prominent figures like Jim Browne, who served on the committee for twenty years and chaired it for ten, was instrumental in shaping campus computing policy and guiding the university towards distributed, interactive networks. This is confirmed by a report from Charles Warlick, the Computation Center director, dated October 10, 1973. Video terminals are referred to as console terminals here, and the expansion of interactive service refers to the acquisition of the DEC-10. This is essentially a roadmap for bringing in video terminals and the DEC-10 beginning in 1974 and 1975. 

\begin{quotation}
The Board of Regents has allocated one million dollars for the further development of instructional computing during the next twelve months. The University-wide Faculty Computer Committee requested the funds and is responsible for their allocation. Approximately one-quarter of the amount will go into console terminals in new teaching centers to be established in the various UT Austin colleges. About half will go into the expansion of the interactive service, with particular attention to the support of students interacting with instructional modules. The balance will go into remote batch computer terminals and equipment for special teaching projects.
\end{quotation}

San Antonio’s Datapoint, originally the Computer Terminal Corporation, is a fascinating story of computing and the birth of video terminals in the late sixties. It was founded in 1968 by two engineers, Phil Ray and Gus Roche. Ray was a Texan and graduated from UT in 1957. He began his career at Texas Instruments, where he contributed to missile telemetry and electronic speech analysis projects, before advancing to roles in aerospace telemetry at International Data Systems and as a senior staff engineer at General Dynamics' Dynatronics division, supporting NASA's Saturn V rocket and Lunar Orbiter programs. In 1967, while working on NASA projects in Florida, Ray partnered with colleague Gus Roche to explore entrepreneurial opportunities amid anticipated space program cutbacks, and events soon took a surprising turn.

\begin{quotation}
Roche and Ray also turned to Bob McClure, then working as a computer consultant in Dallas while teaching as well. Although younger than Ray, he had been a teaching assistant at the University of Texas when Ray was there. Ray had not only been in one of his classes, but the two had worked on an electromechanical demonstration computer, built from relays scavenged from pinball machines, that played a counting game called Nim. “I was sitting in my office in 1968 when I got a call from Charlie Skelton about consulting,” he recalled. Skelton was another Texan and a mutual friend of McClure and Ray, who was at that time involved with Ray and Roche in their effort to start a company. “He and Phil came over and we had a discussion about computer terminals. One of them said they had raised money to start a new firm to build computer terminals. They had read an article in Businessweek saying that computer terminals would be the next big thing but that neither of them had ever seen a computer terminal. By that they were joking, of course. I said I had been thinking about the topic, as I had been doing work in computers and knew that the availability of computer terminals was thin. Nearly all input was done with Model 33 Teletypes at 110 baud. So I said that they should build a glass Teletype. But I told them to make sure that they emulated the protocol of the Teletype exactly, so that no one at the computer end has to change any software, since that would be a hang-up to getting it accepted. There was a CRT terminal already on the market that was not precisely compatible and therefore had a lot of problems. There were some expensive machines available, but I assumed they could make one for less since there was really not that much in one,” McClure recalled.
\end{quotation}

Following through on that initial advice, Ray and Roche went on to create the Datapoint 3300, released in 1969. The Datapoint 3300 was designed, named, and marketed as a silent, electronic glass teletype replacement for the loud, mechanical Teletype Model 33. Introduced in 1963, the Teletype Model 33, often referred to as the ASR-33 for its Automatic Send and Receive variant, was a massively popular electromechanical teleprinter. Because it was relatively inexpensive at around \$1,000, and was one of the first machines to utilize the newly standardized 7-bit ASCII code, it became the ubiquitous terminal by the late sixties. It printed on continuous spools of paper at an agonizingly slow 110 baud, about 10 characters per second, and its heavy, motor-driven hammer mechanisms generated intense, continuous noise. To solve the noise and speed bottlenecks of machines like the ASR-33, Ray and Roche developed one of the earliest standalone commercial cathode-ray tube displays. The number 3300 in its name was chosen to position the device as a direct, futuristic advance over the Teletype Model 33. Datapoint even declared in its marketing that their new terminal was "100 times better than the 33".

To make the transition from paper to screen as seamless as possible for businesses, the Datapoint 3300 was designed to perfectly emulate the Teletype Model 33. It offered complete interchangeability with standard teletypewriter equipment, meaning operators could simply unplug a loud, mechanical Teletype and plug in a Datapoint 3300 to the exact same mainframe connection without requiring any software changes. Because memory was still expensive, the Datapoint 3300 used a digital shift-register design to store screen data and featured a display format of 72 columns of text, slightly less than the 80 columns that would later become standard.

The Datapoint 3300 was an immediate hit when it debuted. In fact, demand so vastly outpaced Datapoint's initial manufacturing capabilities that the company had to temporarily outsource the production of the terminal's sleek, streamlined casings to a local San Antonio motorcycle helmet manufacturer. The terminal proved so successful that major computer vendors purchased them to rebadge and sell alongside their own systems. DEC sold it as the DEC VT06, and Hewlett-Packard sold it as the HP 2600A. The explosive success of the Datapoint 3300 provided the capital and momentum to begin developing a smarter successor terminal with its own internal processing power, a project that would shortly become the revolutionary Datapoint 2200.

The 1970 Datapoint 2200 is considered by some to be the true birth of the personal computer. Though marketed as a terminal, it was actually a programmable desktop system with its own internal memory, keyboard, and screen. It was capable of playing computer games, standalone with no external connection. Its digital architecture and instruction set was reproduced in Intel’s 8008 microprocessor. The Datapoint 2200 is the ancestor of the x86 standard, and every x86 chip instruction set contains 1970 Datapoint instructions and is in a sense backward compatible with the 2200. The initial design relied on a dense, multi-board processor built from roughly a hundred discrete transistor-transistor logic chips. Seeking to reduce the machine's size, heat, and manufacturing costs, in 1969 Datapoint contracted Intel and Texas Instruments to consolidate the entire processor board into a single silicon chip. Neither company could initially meet Datapoint's needs. Texas Instruments produced the TMX 1795, but it was notoriously buggy and ultimately abandoned. 

Intel successfully fit the 2200's exact instruction set architecture into a single chip, but missed the delivery deadline, and the resulting chip was deemed too slow. Intel retained the intellectual property rights and released it commercially in 1972 as the Intel 8008, which became the foundation of Intel's 8-bit processor line and directly inspired the Intel 8080 and the highly successful 8086. This lineage established the x86 instruction set architecture that powered the original IBM personal computer and continues to dominate modern computing today. Because the original Datapoint 2200 processed data serially, it was forced to store numbers starting with the lowest bit first to handle mathematical carries. This little-endian format was inherited by the 8008 and remains the standard for modern x86 processors, and was directly opposed to the far more natural big-endian nature of mainframe class hardware. Some have called Datapoint’s course with regard to the Intel chip one of the worst decisions in business history, as it seems probable that they could have at the least received licensing fees for the instruction set intellectual property.

Because the video terminal market was so new and experimental, there were no standards for what the video screen should look like and how it should function. Datapoint decided that it made sense for the shape and visual feel of the screen to match with standard IBM punchcards, since computer users were familiar with punchcards at the time. This is why the Datapoint 2200 screen looks “squashed” to the modern eye. It matches the shape of a punchcard. In fact, the terminal was designed specifically to replace punchcards, just as the Datapoint 3300 had replaced the Teletype Model 33. The concept was that, in the corporate world data entry was often performed by creating decks of punchcards at remote offices and then sending the decks to a central data processing site. With the 2200, data entry would be done as if it were on punchcards, but on the video screen and keyboard of the 2200 rather than on paper using a keypunch. The 2200 would store the data internally, and write it to a standard consumer audio cassette tape. This system of storing computer data on audio cassette tapes would become familiar to early home computer users. The 2200 actually had two cassette tape drives integrated into its top surface. Remote offices sent cassette tapes to the central data processing site, rather than punchcard decks. This system turned out to be highly effective in practice, and made Datapoint rich by the mid and late seventies.

The evolution of the computer terminal represented a hardware bridge from the mainframe era to the personal computer era. They were both the ultimate refinement of how humans interacted with centralized mainframes, and a catalyst that helped launch the evolution of personal computers. In the earliest days of computing, users interacted with mainframes via front-panel switches, punched cards, and paper tape in a slow, asynchronous batch-processing workflow. The introduction of time-sharing systems allowed multiple users to interact with a single mainframe simultaneously, which drove the adoption of electromechanical teleprinters like the Teletype Model 33. While revolutionary for providing real-time access, these teleprinters were too loud, slow, and wasteful. Video terminals or glass teletypes pushed the mainframe era to its peak interactive potential, allowing users to view and manipulate data at speeds up to 19,200 bits per second. Even with fast screens, remote mainframe connections were often choked by slow, low-bandwidth telephone modems. To reduce the load on the central mainframe and minimize the data sent over these slow lines, manufacturers began putting microprocessors into terminals so they could handle basic editing and data validation locally. These devices were dubbed intelligent terminals.

The most pivotal of these was the 2200. It was marketed to businesses as a highly flexible programmable terminal that could emulate various other terminals by loading software from its built-in dual cassette drives. To make it friendly for office use, it was housed in a compact desktop case designed to mimic the footprint of an IBM Selectric typewriter. In reality, Datapoint had unwittingly designed the first commercial personal computer. The transition from terminal to personal computer became undeniable when customers started ignoring the mainframe entirely. In 1971, one of the 2200's designers, Victor Poor, visited a Pillsbury chicken-processing plant in Arkansas to observe the terminal in action. When he asked the local managers what mainframe they were hooked up to, they informed him that the Datapoint 2200 was running their payroll and inventory locally. It wasn't connected to a host computer at all. The terminal had become a standalone computer.

In 1978 the DEC VT100  jumped from the rigid, hardwired circuitry of earlier terminals to a flexible, microprocessor-driven architecture powered by an internal Intel 8080 microprocessor, so was effectively a direct descendant of the 2200. By using a microprocessor, the VT100's behavior was governed by software rather than physical hardware in the form of complex arrays of discrete logic chips. The older transistor-transistor logic architecture gave way to the microprocessor architecture that would drive the personal computer market. In some sense, the VT100 was trending rapidly towards a personal computer. Compared to the 2200, it was locked up by design for marketing reasons, but because of the Intel 8080 microprocessor the trend was clear. If the VT100 could be jail broken, it could make a decent personal computer.

The processing power of the 8080 enabled the VT100 to be one of the first commercial terminals to parse the newly emerging ANSI X3.64 standard for command codes. Rather than forcing a mainframe to resend an entire screen of text over a slow modem line just to update a single word, the host computer could send highly efficient escape sequences. The VT100’s microprocessor would interpret these codes locally to position the cursor at specific coordinates, clear lines, or insert characters, significantly reducing the bandwidth required for text editing and interactive applications. Because its operations were software-defined, it was able to include a fallback mode that could seamlessly emulate the older, proprietary escape sequences of the VT52. This ensured that existing software and mainframe applications would continue to run flawlessly without modification, while still granting users access to the new ANSI features when needed. 

Previous terminals required operators to manually toggle physical switches or jumper wires to change communication settings. The VT100 used its microprocessor to generate an interactive, on-screen setup menu where users could configure their terminal. The processor then saved these preferences directly into the terminal's non-volatile memory, ensuring settings were retained even when the machine was powered off. The microprocessor allowed the terminal to independently manage dynamic graphic renditions, meaning it could instantly display text with bolding, blinking, underlining, reverse video, or double-sized characters. It also featured a special box-drawing character set to locally render on-screen forms. It introduced smooth scrolling, which moved lines of text gradually up or down the screen to make reading much easier, replacing the sudden, jarring jumps that plagued earlier terminals. Because the microprocessor could manage incoming data and screen refreshes much more efficiently, it supported maximum serial communication speeds of 19,200 baud. This effectively doubled the maximum data rate of its predecessor, the VT52, making full-screen interactive computing much faster and more responsive.

It took the October 1973 UT Computation Center roadmap, the financial backing of the Board of Regents, and the arrival of industry-standard smart terminals like the DEC VT52 around 1975 to make interactive computing an institutional reality. The interface between the user and the mainframe evolved from the electromechanical Teletype and DECwriter to industry standard video terminals like the VT52 and VT100, alongside more specialized hardware like Datapoint systems and Tektronix vector graphics terminals. Despite these advances, the operational reality remained constrained by low-bandwidth serial line connections. When it was possible at all, dial-in access via early modems was limited to 300 or 1200 baud making remote interaction painfully slow. In general, because the terminals and network connections were still primitive, programmers expected to print massive fanfold paper listings to conduct manual code reviews. The days of punchcard decks may have been over, but teetering and dusty stacks of greenbar paper output was still very much a daily reality.

\section{Simulation And AI}
In the fifties, once assembly had become conventional and Fortran had followed in its footsteps, the two were used together side-by-side as a historically distinctive pair. For pure floating-point modeling of physical systems, it was standard to write inner loops in assembly and outer loops or one-time code in Fortran, forming a kind of family that defined an era and a group of users associated with numerical simulation. Understanding the assembly and Fortran symbiosis provides an insight into the early days of computing and forms an obvious archaeological layer in computing history. The Decwar project at UT was certainly one of the very last to be approached from the classic assembly and Fortran world-view. Contemporary projects such as Utopia, created by Bob Schutz and his research group, were written primarily in Fortran and rarely or never explicitly used assembly because the necessary assembly-level work was taken care of in standardized, portable numerical libraries. In some sense, Decwar and Utopia are contrasting examples in the long history of assembly and Fortran, and in the story of modeling and simulation at UT in the seventies. For that era, both were effectively large scale simulations, one with a satellite traversing the Earth’s gravity field and interacting with a network of ground tracking stations, the other with up to eighteen people battling across a galaxy of stars, planets, novas, and black holes.

By the seventies, a very different spirit was taking root in parallel, and from early on it was identified with artificial intelligence in the broad sense of doing things more in the way that humans do, for example using human language or reasoning about and manipulating items in the real-world environment. This was a completely different world from the kinds of modeling and simulation associated with the early days of computing and the assembly and Fortran symbiosis. New tools and projects entered the stage. An example is the german to english machine translation Metal project in the Linguistics Research Center, in which older Fortran code on the CDC 6600 was directly replaced by Lisp code on the DEC-10.

Before exploring the new type of projects appearing by the late seventies, the stage can be set by looking at the situation in early 1975 when the DEC-10 was newly installed. The PDP-10 processor at the heart of the system was a 36-bit word-addressable machine, a design with advantages for both character manipulation and high-precision mathematics. The 36-bit word was the fundamental unit of storage, but the processor was capable of handling 18-bit halfwords, which were used extensively for pointers and addressing, since the memory address space was 18-bit. In the UT environment, this memory was managed through a sophisticated paging system, eliminating the need for contiguous memory allocation and reducing the overhead associated with checkerboarding in physical core memory. In addition to single and double word floating-point numbers, a single word could represent five 7-bit characters, with one bit unused, or six 6-bit characters. The two options were usually identified as Ascii and Sixbit, respectively. The first sixteen words of memory were actually fast registers within the processor hardware, though they remained addressable as memory locations. Assembly programmers frequently exploited this by loading small, frequently executed loops directly into these registers to serve as an impromptu instruction cache, a technique notably used in early Lisp implementations on the machine. 

For systems programmers and technical staff, assembly was the essential instrument for system optimization and infrastructure maintenance. As a two-pass symbolic assembler, Macro-10 provided a direct and elegant interface to the hundreds of hard-wired instructions of the processor, facilitating a degree of optimization unachievable by higher-level compiled languages. The language's rich set of pseudo-operations empowered university staff to manage intricate memory layouts, reserve shared memory segments, and generate highly modular code for local utilities. A defining characteristic of assembly programming was the extensive use of unimplemented user operations. Here unimplemented meant simply that they were not implemented as hard-wired instructions, but as user software instead. These allowed programmers to interact intimately with the operating system monitor to directly command file input and output operations, allocate disk space, govern inter-process communication, and enforce local administrative policies such as custom disk quotas, without destabilizing the operating system's kernel. 

While assembly was the language of the machine, Fortran was the primary engine for higher-level engineering and scientific research, including demanding applications in aerospace, astronomy, and operations research. Multiple Fortran compilers were available, but the default Fortran-10 was the standard choice for performance-critical tasks and is simply assumed as synonymous with Fortran here. It featured one of the most advanced optimizers of its era, capable of advanced register allocation and common subexpression elimination to allow Fortran programs to run at speeds that often rivaled hand-tuned assembly code. And, critically, Fortran not only supported but emphasized the concept of shared memory, referred to as the shared upper-segment. The upper reaches of the memory space could easily be shared between multiple jobs all accessing the same data values, with Fortran and assembly cooperating to provide system level locking and unlocking to prevent collisions between jobs. Despite its extensions, it maintained a high degree of compatibility with international Fortran standards, ensuring that research code was reusable and could be shared with other institutions using different architectures.

An entirely different direction was opened with the Certifiable Minicomputer Project, established in 1975 by Don Good to prove that mathematical verification could replace empirical software testing for mission-critical and highly secure systems. The project developed Gypsy, a high-level programming language that combined executable statements with formal mathematical specifications in a single syntax. Gypsy was designed as a high-level programming language and formal specification environment derived from Pascal. While the language itself was a Pascal derivative, the systems used to process and verify it were built using other languages. The surrounding tools, such as the System for Incremental Development used within the Gypsy Verification Environment, were written entirely in Lisp. These tools also included a Verification Condition Generator and an interactive theorem prover. 

The CMP was deeply involved with UT Austin’s connection to the Arpanet in 1977, through a landmark software verification project known as the Encrypted Packet Interface. The Arpanet required a secure method for communicating computers to encrypt messages as they entered the network and decrypt them as they returned. To address this, CMP researchers developed a 4,200 to 5,000-line software program written in Gypsy and engineered to precisely duplicate the network encryption and decryption functions of a proprietary hardware black-box that had been developed by Bolt Beranek and Newman. The CMP used the Gypsy Verification Environment to construct a mechanical proof of its behavior, and the Gypsy software achieved perfect interoperability with the real-world BBN hardware. This was a milestone for the field of formal verification, as it established that mechanically proven software could successfully and securely replace specialized, physical hardware components in critical networks like the Arpanet. Another milestone was the Message Flow Modulator, a security filter program designed for the Naval Electronic Systems Command. The MFM was fully formally specified and proved mechanically. Tested independently at the Patuxent River Naval Air Test Center, the verified MFM passed every test on the first attempt without a single code modification, becoming the first mechanically proven software system successfully deployed in real-world military security operations.

Other Lisp based projects were associated with the Linguistics Research Center. The LRC was founded in 1961 by Winfred Lehmann and was in the vanguard of early computational humanities, using mainframes to transcribe, analyze, and publish ancient texts. Heavily funded by defense contracts, its early machine translation efforts focused on german to english translation using a monolithic Fortran-based system on the CDC 6600. This fit well with the extensive german heritage in Texas, especially in central Texas near Austin. When federal funding for machine translation collapsed following the critical 1966 ALPAC report, the LRC's systems fell into disrepair. In 1978, the project was resurrected by the german firm Siemens to meet a corporate need for high-speed translation. This revival retained the earlier Metal project name but was really a major change from the earlier work. When Jonathan Slocum took over in 1979, he scrapped the unusable Fortran code and rebuilt the Metal system in Lisp on the DEC-10. Slocum redirected machine translation by shifting the bottleneck from programmers to linguists. He built specialized graphical user interfaces, real-time rule validators, and an automated tool that allowed non-programming linguists to build, test, and debug massive lexicons and grammatical rules interactively. A full production prototype was delivered to Siemens in 1984, laying the groundwork for early commercial translation tools.

Other computational linguistics projects at the LRC grew around the DEC-10. Lauri Karttunen’s research led to Kimmo, the first operational, language-independent computational parser, implemented in Lisp. Frances Karttunen directed the Mesoamerican Languages Project at the LRC, where her work focused on Nahuatl and Yucatec Maya. The Karttunens collaborated on Frances's landmark Analytical Dictionary of Nahuatl. This complex dictionary systematically documented crucial but historically ignored features like vowel length and the glottal stop, with the direct support of Lauri’s database algorithms. Their research demonstrated how the computational database methods developed during this era could effectively preserve and document endangered indigenous languages.

The Cognitive Science program at UT combined computer science, psychology, linguistics, philosophy, and neuroscience to study how both biological and artificial agents represent knowledge, process complex information, and navigate linguistic context. Robert Simmons began the program in 1968 after moving from the RAND spinoff System Development Corporation in Santa Monica, another direct connection from UT back to RAND, SAGE, and the Q-32. While serving as the head of the Language Processing Research Program at SDC, Simmons directed the Synthex and Protosynthex projects, which aimed to synthesize human language behavior, read texts, and generate English answers to natural language questions. This research relied heavily on Lisp, and SDC became a major hub for Lisp development, because navigating complex semantic networks required dynamic memory allocation and advanced list-processing capabilities.

The physical hardware hosting this research, the AN/FSQ-32, was a fascinating artifact of the Cold War. Originally engineered as a prototype for the military's SAGE air defense system, its deployment was canceled, allowing ARPA to repurpose the machine at SDC as an open-ended laboratory for timesharing, networking, and symbolic processing. This unique environment produced groundbreaking Lisp tools. The Q-32 Lisp 1.5 port was pioneering because it discarded the traditional interpreter-first execution model. It compiled all user code directly into machine instructions prior to execution to maximize speed. SDC and its collaborators even attempted to build Lisp 2, an ambitious language designed to combine Lisp's list-processing capabilities with the structured, algebraic syntax of Algol. Ironically, the Q-32's hardware limitations, specifically its absolute ceiling of 48,000 words of core memory, proved fatal for the highly complex Lisp 2 system, which routinely exceeded the machine's capacity and led to the project's cancellation in 1967.

To understand exactly why Lisp 2 exhausted the Q-32's memory space, it helps to look at everything the development team was trying to pack into it. Beyond adding the Algol syntax, the system featured a complex meta-compiler, dynamic arrays, multiple data types, and a highly advanced, retargetable compiler designed to eventually port the language to machines like the IBM 360 and DEC PDP-6, direct ancestor of the DEC-10. To force this massive environment into the Q-32's cramped core memory, developers had to engineer an intricate, application-level memory swapper. The system's execution pipeline became heavily bogged down in memory management tasks. If a required function wasn't in active core memory, the system had to page the compiled binary from an external storage file. When memory was fully occupied, the manager would initiate an immediate compaction cycle, shuffling the in-memory binary code to assemble contiguous free space. If compaction failed to yield enough room, the system would selectively excise non-essential routines or trigger a full garbage collection pass to reclaim list storage before it could finally load the new function.

This constant disk paging and memory shuffling introduced severe latency. The overhead choked the system's performance and severely degraded developer productivity, making it nearly impossible to rapidly iterate on the language's syntax translator and optimization routines. These technical roadblocks were compounded by rising costs and collaborative friction among the project's institutional partners SDC, MIT, and Stanford. When the project was ultimately canceled in 1967, the AI community abandoned Lisp 2 and migrated back to Lisp 1.5, finding a more suitable home on the PDP-6 and its descendant PDP-10 architectures. Interestingly, LISP creator John McCarthy later lamented the cancellation, noting that "much more money has since been spent to develop LISPs with fewer features". However, at the time, the developers could not have foreseen that the PDP-10 architecture, which had a much larger address space, would soon emerge to solve the very hardware limitations that doomed Lisp 2.

When Simmons moved to UT Austin in 1968 he was able to escape these specific hardware constraints and expand his work. To support his growing research group, UT bypassed the need for remote teletype connections back to SDC by implementing its own native version of Lisp 1.5 on the CDC 6600. This localized computing power democratized access to Lisp, enabling Simmons and a new generation of graduate students to run the large-scale symbolic experiments necessary to pursue his ultimate dream of enabling a "conversation with a book", where a computer system could read expository text, map the underlying knowledge into semantic networks, and generate natural language answers to users' questions. 

The decision to build a native Lisp system at UT was driven by the sheer unreliability of remote computing in the late sixties. Relying on remote teletype connections back to SDC in Santa Monica meant dealing with manual dialing, fragile acoustic couplers, and unstable long-distance phone lines. Because Lisp's parenthetical syntax is highly sensitive, a single corrupted character from line noise would cause the parser to reject the input or trigger a runtime crash. To overcome this, the university utilized its CDC 6600 and 131,072 words of magnetic core storage to build a native Lisp implementation written directly in CDC Compass assembly language. This system, known as UT-Lisp, introduced a radical and highly unique architectural innovation referred to as the three-pointer cons cell.

Because the CDC 6600 used a 60-bit word and an 18-bit address space, a traditional Lisp cell containing only a CAR and a CDR pointer, with 18 bits each, would leave 24 bits empty and wasted. UT-Lisp's designers partitioned the 60-bit word into three distinct 18-bit fields. The usual CAR and CDR fields, and a special field CSR. To navigate these unique three-pointer cells, the system supported compositional accessors of extraordinary depth. Programmers could chain up to eight operations together, which the system's handler would dynamically parse and execute at runtime. To manage the memory constraints of the university's multi-user timesharing environment, Mabry Tyson also developed a specialized virtual memory system for UT-Lisp that paged individual Lisp functions in and out of core memory dynamically.

This unique three-pointer architecture proved perfectly suited for Simmons' research into cognitive science and computational linguistics. His group focused on representing meaning through semantic networks, which map concepts as nodes and relationships, using labeled directed arcs. Conceptual triples with a subject, relationship, and object mapped naturally onto UT-Lisp's three-part CAR, CDR, CSR pointer structures. Using this highly customized local environment, Simmons and his research group were able to bypass the physical memory bottlenecks that had previously killed ambitious projects like Lisp 2. Leveraging UT-Lisp, they successfully migrated Protosynthex from the Q-32 to the CDC 6600 to parse English strings, and completely rewrote the Linguistic Research Center’s Fortran based Metal machine translation system into Lisp for real-time translation.

Another UT-Lisp project was Lawaly, a general-purpose robot planner and problem-solving system developed in the early seventies by Siklossy and Dreussi. Programmed in UT-Lisp and run interpretively on the CDC 6600, Lawaly was designed to emulate how a human being might write a specialized problem-solving program for a given environment. Instead of relying on hard-coded routines, Lawaly was fed an axiomatic description of a specific environment and its capabilities within it. From these axioms, it dynamically generated its own procedures by creating a specific Lisp executable expression for each operator. To solve complex problems, it relied on a hierarchy of subtasks. If it failed to achieve a group of subtasks, it would systematically test permutations of those tasks and reuse parts of the failed paths to avoid dead ends, minimizing the need for destructive backtracking.

Lawaly was highly adaptable and successfully applied to more than twenty distinct simulated worlds. In some of the more complex environments, it operated as a robot-janitor. Its generated procedures allowed it to navigate worlds where it had to sweep floors, empty trash baskets into dustbins, fill pails at faucets, dry terraria with hot winds, turn lights on and off, and climb onto boxes. Some of these worlds were massive for the era, containing over 500 distinct static and dynamic predicates. Lawaly was extensively evaluated against Strips, the famous theorem-proving robot planner developed at the Stanford Research Institute. While Strips often struggled with complex problems unless it was provided with macro-operators that generalized sequences of previous tasks into single steps, Lawaly routinely solved tasks requiring several hundred steps without needing any prior learning. Because of its hierarchical approach, it's time to find a solution grew linearly with the number of steps required, rather than exponentially. It exhibited phenomenal execution efficiency, averaging just a second or two per node in its search tree. Even accounting for the fact that the CDC 6600 was roughly eight times faster than the PDP-10 that Strips ran on, Lawaly’s overall solution times remained vastly superior. Despite its power, Lawaly had a few specific limitations. It lacked maze-running capabilities, meaning it had no inherent sense of physical direction when navigating rooms. Its strict focus on hierarchical subtasks occasionally led to amusing failures, a flaw its developers characterized as stubbornness.

In 1972, Simmons co-authored the paper "Generating English Discourse from Semantic Networks" with his student Jonathan Slocum. A few years later, in 1979, Slocum was hired to take over the LRC's faltering Metal machine translation project because of his proficiency in natural language processing. To save the project, Slocum completely scrapped the old CDC 6600 Fortran code and rebuilt the entire translation engine from scratch in Lisp, creating a system that finally allowed real-time linguistic rule testing. Simmons also co-authored the 1974 paper "Modeling Dictionary Data" with Robert Amsler. Amsler served as a critical computational bridge between different language projects, providing technical and computational support for Frances Karttunen's Mesoamerican Languages Project at the LRC. Later in his career, Simmons became interested in logic programming and natural language processing based on logic programs. In the context of AI history, early logic programming, such as the Micro Planner language used to build the famous natural language AI system Shrdlu, was implemented directly within Lisp.

While the CDC 6600 was busy hosting ambitious symbolic artificial intelligence projects like Lawaly and the Metal translation system in UT-Lisp, its Fortran environment was simultaneously serving as an incubator for early computer gaming culture. As will be seen, the ancestor of Decwar’s computer-controlled Romulan was born on the CDC 6600, and would become a highly effective demonstration of the possibilities for AI. For many Decwar players, the Romulan was a far more memorable AI experience than the contemporaneous Eliza or Doctor programs.

\section{TECO-124, MACRO-10, FORTRAN-10, DECWAR}

Decwar originated as a Fortran single-player Star Trek game, created by Dave Matuszek, Paul Reynolds, and Rich Cohen on the CDC 6600 in 1973 and 1974. This initial Fortran version was itself an effort to move beyond an earlier resource-intensive Basic version created in 1973 by Grady Hicks and Jim Korp, and inspired by the original 1971 Mayfield Star Trek game championed by Dave Ahl in his highly influential 1973 book 101 BASIC Computer Games. The Basic version was slow and problematic for the Computation Center staff, who reportedly disliked it for loading down the system. The Fortran version immediately bypassed the issues associated with Basic and won support from the staff. Over the next few years, Robert Schneider and others extended the Fortran code to a two-player version on the 6600 named War, with two players taking turns at a single teletype. The name refers to the old children's game Paper War, often called Pencil Wars or Paper Battles, a classic two-player game where opponents draw troops on a folded sheet of paper and take turns attacking using blind-folding marks or physical pencil-flicking dexterity. 

When the time came to move to the DEC-10 after it arrived in 1975, Robert Schneider was involved with transferring the single and two player code from the CDC, and at some point student-developers Bob Hysick and Jeff Potter began to take over the evolutionary process, fundamentally altering the overall trajectory to take advantage of the DEC-10's architecture and operating system. Computation Center experts actively shared their knowledge and guidance, in particular Clive Dawson, Tommy Loomis, and Rick Watson. The project name evolved naturally from War to Decwar, aptly recognizing both the legacy of the CDC 6600 two-player game and the 1962 Spacewar game from MIT, and the code was rewritten as a tight pairing of high-level Fortran and low-level assembly language. The ambition was to dramatically expand the scope and simulate a galaxy filled with players, a goal that inherently meant pushing the operational limits of the system. To achieve a persistent, concurrent environment where multiple users could interact with a shared game state, they needed to leverage the DEC-10’s full capabilities, from the operating system to the assembler and Fortran compiler. The development of Decwar was the direct result of extending the classic mainframe-era Fortran and assembly symbiosis, and the lineage of the game traces a direct path through the Computing Center’s history over the course of the sixties and seventies, marking a fitting culmination for the era. 

The DEC-10's architecture and operating system provided a critical feature that made massive multiplayer gaming possible, a shared high segment memory to store the state of the galaxy. This foundational support for shared memory and inter-process synchronization via locking were crucial. It was obvious that scaling the game to support numerous simultaneous players required robust mechanisms for shared memory access and a system-level locking mechanism to prevent concurrent data collisions within the magnetic core memory. Instead of a simple two-player duel, the game now allowed up to eighteen people to log in simultaneously. Each player's ship was controlled by a completely separate, independent job, and all eighteen jobs interacted in real-time by reading and writing to a single shared block of core memory. To manage the chaotic concurrency of eighteen players attacking and communicating at once, the developers had to rely heavily on assembly to interact with the operating system's native locking mechanisms, preventing race conditions when multiple jobs tried to alter the same memory block simultaneously. 

To force Fortran and assembly to share this delicate memory map, the Decwar developers built a sophisticated, custom toolchain that serves as a testament to the hacker ingenuity of the era. Long before modern build systems existed, they depended on executable macros or scripts within the Teco environment to auto-generate code. It became so essential to the game's operation that the UT version, Teco-124, was bundled directly onto Decwar software distribution tapes along with the game. Before compiling and linking the game, these Teco macros would parse the master assembly files and programmatically generate Fortran include files. 

Teco was originally created in 1962 by Dan Murphy for use on the PDP-1 computer at MIT. This was the same year that Spacewar, the ancestor and namesake of Decwar, was created on the same machine. At the time, it stood for Tape Editor and Corrector because the PDP-1 lacked hard disks or magnetic tape. Program source code was stored entirely on continuous strips of punched paper tape. To avoid the tedious physical manipulation of paper tape used by early programmers, Murphy designed Teco to allow a user to feed a source tape and a "correction tape" into the machine, which would then rapidly punch out a newly edited tape. As storage media evolved, the acronym was updated to Text Editor and Corrector, but its highly cryptic, character-oriented nature remained. Teco did not use numbered lines. It maintained a pointer between characters, allowing editing operations to span across line boundaries. Because it lacked a formal syntax, every character typed was an immediate, imperative execution command, and typographical errors could easily corrupt a file. This earned Teco the satirical acronym YAFIYGI You Asked For It, You Got It, and inspired a popular game where programmers would type their names as command sequences just to see what the interpreter would destroy.

Clive Dawson developed Teco-124 on the UT DEC-10 with video terminal display logic integrated directly into the compiled binary. This gave programmers using video terminals native cursor movement and full-screen visual editing capabilities. By the mid-seventies, Digital distributed a standard version Teco-24 with the DEC-10. Because this version was created during the teletype era, it lacked the ability to process the escape sequences required for real-time visual screen updates on smart video terminals. Teco-124 was UT’s answer. Dawson gives an example of its use in the following quote.

\begin{quotation}
Across the hall at the Linguistics Research Center was Dr. Helen Jo Hewitt, aka “HJ”.  She was a BIG fan of Teco, and was one of my “power” users.  She was one of the early “font hackers”, and I had developed a bunch of Teco macros for her.  She would come across the hall and pose an interesting problem, and a few minutes, or hours, or days(!) later, I would present her with another Teco macro for her growing collection. These tools would enable her, for example, to use a meta-language of her own creation to insert diacritics into a document, and then apply the appropriate macro to convert the meta-language into actual escape-sequence commands for the robotic typewriter to backspace, make micro-adjustments to the carriage, and type the appropriate symbol(s) to produce the desired diacritic at j-u-s-t the right position. Later on, as her daisy-wheel collection of different type-fonts grew, a macro would be used to type a partially-filled page of text, then command the typewriter to roll the page back to the beginning and pause. She could then switch in, say, the Greek daisy-wheel, after which the job would continue with the typewriter skipping directly to each blank spot where a Greek character needed to be typed.    
\end{quotation}

Teco's legacy is ultimately defined by what it evolved into. As video terminals became available, the MIT AI Lab’s hackers added a display-editing mode to Teco known as control-R to update the on-screen layout in real-time. This matched closely with Dawson’s work with UT’s Teco-124. In parallel, Teco was evolving into Emacs. In 1974, coders expanded MIT’s display-editing mode by adding a macro execution feature, allowing users to bind keystrokes to Teco programs to automate common tasks. Users quickly accumulated a diverse collection of custom macros. In 1976, these macro sets merged into a single, standardized, extensible system called Emacs, short for Editing Macros. Though Emacs began as a collection of Teco macros, it eventually outgrew its host. It was rewritten into a standalone program and finally transformed into GNU Emacs in 1984, trading Teco for a true Lisp runtime. The legend soon spread that with Emacs, Teco was achieving sentience and signalling the true dawn of AI. It all started innocently enough. You opened Emacs to edit a simple .txt file, and suddenly it was running your email, managing your calendar, playing Tetris, controlling your smart lights, and hosting its own web server. A text editor doesn't just grow a built-in Lisp interpreter and a terminal emulator without having higher ambitions.

For Decwar, maintaining strict alignment between symbolic variables in the high-level Fortran routines and the absolute memory offsets manipulated by the Macro-10 assembly code was a significant challenge. To manage this, the developers constructed an automated code-generation toolchain based on Teco macro scripts. These scripts scanned the master assembly files and automatically extracted parameters to generate Fortran include files. These tools generated Hiseg.for and Lowseg.for, which defined the Fortran common blocks that mapped variables directly into shared high segment memory and the private low segment memory of individual player jobs. By generating Extern.for from the assembly file Msg.mac, developers were even able to force user-facing text strings entirely into the shared high segment, minimizing the memory footprint of player jobs.

Teco was used at UT from 1975 for working with the Tops-10 operating system. Video terminals were rapidly becoming common, and Dawson’s Teco-124 took advantage of Teco’s inherent flexibility to make the user experience more interactive. At the same time, the operating system was opening up the possibility of real-time multi-user interaction. The combination of video terminals and multi-user interactivity was liberating. All of this was supported by the operating system through software additions to the instruction set known as unimplemented user operations. These were effectively kernel-level system calls that, for example, allowed for the atomic locking and unlocking of memory segments, preventing data race conditions. This locking mechanism enforced sequential access, ensuring that only one user could acquire the lock to modify a specific word in the shared memory at a time, although reading remained open to all. Because Fortran lacked the native syntax to invoke these sophisticated system-level primitives, the hybrid programming model was adopted, marrying the algorithmic structure of Fortran with the hardware-level precision of assembly. 

Macro-10’s name suggests that it was in some sense more than just an assembly language, and its resonance with the recurring theme of macro programming in the context of Teco, Emacs, and Lisp is suggestive. The name was a deliberate reflection of a broader theme within the system’s landscape, the reliance on extensive macro programming to achieve open-ended extensibility. The boundaries between programming languages, text editors, and runtime environments were remarkably fluid. Just as Teco allowed users to write complex editing macros, which eventually spawned Emacs, and Lisp relied heavily on macros to extend its syntax, Macro-10 treated the assembler itself as a highly programmable macro processor. This was a welcoming language that formed an essential part of the legendarily rich DEC-10 environment. 

Normal users saw this assembler as an everyday tool, and it represented simply a standard approach for taking full advantage of the system. There was nothing esoteric or intimidating about it. Rather, it was a democratized, highly approachable development environment that served as the lingua franca of the interactive computing culture. Pieces of Macro-10 code were written as freely as pieces of Python or Javascript code are written today, and often for similar reasons. In fact, it’s safe to say that any real programmer expected to use Macro-10. It was simply part of the nature of being a programmer. Because of this accessibility, Macro-10 was considered an approachable starting point rather than an intimidating barrier. Introductory books were written for complete beginners, notably Singer’s 1978 “Introduction to DECsystem-10 Assembler Language Programming” and Gorin’s 1985 “Introduction to DECsystem-10 Assembly Language Programming“. Students routinely used Macro-10 to develop everyday utilities, text manipulators, and early interactive games. It challenged the assumption that low-level assembly language is inherently intimidating, and proved that an assembly language could function not just as a rigid machine translator, but as a highly expressive, human-centric development environment. Here is Singer from his book’s introduction.

\begin{quotation}
In our opinion, however, the most useful function served by a knowledge of assembler language programming is to give the user a much closer awareness of how the computer works, as well as inestimably greater control over its workings, than is feasible with a higher level language. In our experience, the higher level language user who is familiar with assembler language is a more efficient, even a happier, programmer than the one who is not. Every computer facility supplies booklets explaining the login procedure, by which the user gains access to the computer. The novice is then too often left facing across a chasm, beyond which, hopelessly out of reach, lie the manufacturers' manuals and many superb texts on the theory and practice of programming. This book is intended to serve as a bridge across that chasm. It is suitable for use by the higher level language user who would like to learn assembler language; but also, we would like to stress, by the complete beginner with no knowledge whatsoever of computers. The notion that assembler language programming is esoteric and inherently difficult is, in our experience, very much mistaken. On the contrary, for many people it seems to be the natural way to start off with computers.
\end{quotation}

Besides allowing users to transparently override the machine's hardware instructions and transform Macro-10 into a fully customized, domain-specific programming environment, what truly made Macro-10 feel like a modern, interactive language was its integration with the Ddt dynamic debugging technique symbolic debugger. In an era where most computing involved a detached, frustrating loop of writing, compiling, running, crashing, and restarting from scratch, Ddt treated the active, running memory image of the program as the primary workspace. Because the symbolic definitions were preserved in memory, programmers could use Ddt to disassemble machine code back into its symbolic Macro-10 mnemonics in real time. If a bug was found, the programmer did not need to restart. They could type new Macro-10 statements directly over the disassembled instructions, patch the memory image using dynamic buffer regions, and seamlessly resume execution. This conversational, real-time workflow made writing assembly code feel like working in a high-level interpreted environment. Perhaps the most famous testament to Macro-10's utility as a flexible development engine was the creation of Altair Basic by Bill Gates and Paul Allen, in the first days of Microsoft circa 1975. Lacking the interfaces to develop directly on early microcomputers, they used Macro-10 on a Harvard PDP-10 to write an Intel 8080 emulator. By defining a collection of Macro-10 macros, they mapped Intel 8080 instructions to equivalent PDP-10 instructions and interactively debugged the Basic source code inside Ddt.

In the Macro-10 context, a macro is a user-defined sequence of assembly statements that are expanded and generated in-line, replacing the macro call itself whenever the macro's name is invoked in the program. This enables sophisticated compile-time metaprogramming. The standard syntax is Define macroname(dummylist) <macrobody>. The symbolic macroname is used to call the macro. An optional, comma-separated list of dummy arguments or parameters is enclosed in parentheses, followed by the body sequence of statements to be generated, enclosed in angle brackets < >. When a macro is called, the assembler performs a textual substitution, replacing the dummy arguments in the macro body with the actual arguments passed in the call. It provides several advanced features for handling these arguments. Programmers can assign default values to arguments so that if the argument is omitted during the call, the assembler automatically fills in the default. If a macro contains labels and is called multiple times, it would normally trigger a duplicate label error. By prefixing a dummy argument with a percent sign, it automatically generates a unique, sequential local label for each individual call. The apostrophe character serves as a concatenation operator within the macro body. It joins a passed argument with neighboring text, and the assembler removes the apostrophe during the expansion. By prefixing an argument with a backslash or quote, the programmer can force the assembler to evaluate the mathematical expression and pass its Ascii or Sixbit string representation rather than the raw text.

A central problem within Decwar was how to identify players, and the solution involved the kinds of bit-level manipulation that assembly code inherently excels at. Some of the most common operations within the game involved looking-up the data associated with a particular player, and this was a fundamental bottleneck for overall performance. If any part of the user identification code was slow, then it set the bounds on the overall performance of the simulation. Within the rigidly sequential environment of the 36-bit PDP-10 processor, the solution was a half-word 18-bit identification system, a choice that established a hard ceiling of eighteen players. Each player was assigned a bit within a half-word, bit one for player one, bit two for player two, etc. The same half-word identification system would result in sixteen players on a modern 32-bit processor, and thirty-two players on a modern 64-bit processor. This bit-level system is actually related to the Ascii and Sixbit schemes for identifying text characters. If six or seven bits together could act as a symbol, then single bits could as well, and in some sense more directly and purely.

Because of its unique symbol resolution hierarchy, Macro-10 searches the macro symbol table before checking the hardware op-code table. This allows programmers to define macros that share names with native hardware instructions, effectively overriding the machine's hard-wired operations. This capability allowed programmers to transform the systems assembly language into a fully customized, domain-specific programming environment. The macros support complex logic and iteration, allowing them to function like high-level language constructs. They can iterate over a list of passed sub-arguments, iterate character-by-character over a string, or loop a block of code a specific number of times. Conditional directives allow generating different assembly code based on the parameters passed or the system environment, and macro can be nested up to eighteen levels deep. With these powers, the assembler begins to seem more like a higher-level environment, and this is a source of PDP-10 assembly’s mythic reputation. The low-level and high-level are blended together in a harmonious and artful balance that is an aesthetic pleasure to experience.

At the same time, Fortran-10 was meeting the Macro-10 assembly code halfway, somewhere in the gray area in between the low-level hardware and modern high-level languages. In 1954, when John Backus proposed an "automatic programming" system, the primary challenge was the extreme scarcity and cost of processing cycles and memory. The prevailing sentiment among the era’s elite programmers was that no compiler could ever match the efficiency of a human expert writing in assembly. To prove them wrong, Fortran’s creators focused on the implementation of the first optimizing compiler capable of translating mathematical expressions into IBM 704 assembly code with such precision that the performance gap essentially vanished. Fortran was not merely the first successful high-level language, it was a sophisticated replacement for manual assembly coding. It is a language designed explicitly to offer the transparency and efficiency of assembly code while providing a portable mathematical notation. The IBM 704 hardware directly influenced the language's syntax. The inclusion of features such as subscripted variables for arrays was designed to exploit the 704’s index registers. Even the iconic three-way arithmetic decision statement was a direct high-level representation of the hardware's branching logic based on the accumulator's sign bit. This proximity meant that a Fortran programmer was essentially writing "structured assembly" where each high-level construct had a predictable, one-to-one or one-to-few mapping to machine instructions. 

The philosophy of hardware-software co-design that defined early Fortran remained a driving force well into the sixties and seventies. The desire to build physical architectures specifically tailored to execute Fortran efficiently was famously echoed by supercomputer pioneer Seymour Cray. When a user complained about the CDC 6600's poor performance with the APL programming language, Cray bluntly replied that he wanted a computer that would run Fortran, so he designed and built one to do exactly that. As Fortran transitioned into the DEC-10 era, the mandate to match the efficiency of manual assembly coding persisted. The system's flexible architecture and instruction set presented a unique challenge in this regard, offering so much scope for compiler optimization and bit-fiddling that compilers had to become exceedingly complex just to capture the numerous tricks a hand coder naturally used. To close this performance gap and reduce the need for manual assembly coding, the Fortran-10 compiler was built around a highly sophisticated global optimizer, employing advanced machine-aware techniques. These optimizations yielded object code that executed up to forty percent faster than unoptimized code and required only half the processor time to compile compared to earlier compilers. Its relevance to hardware access and assembly-level performance is rooted in three foundational pillars: a deterministic static memory model, low-level memory mapping constructs such as common and equivalence, and a subroutine linkage convention that defined the early standards for application-level binary interfaces. 

By the mid-sixties, more than forty different Fortran compilers existed, each with minor variations. In 1966 came the first national computing standard, ensuring that Fortran scientific code could be moved between disparate architectures such as the CDC 6600 and DEC-10, while maintaining high-performance access to the underlying hardware. A fundamental reason for this was rigid adherence to static memory allocation. Unlike modern languages that rely on dynamic heap allocation or complex stack frames, Fortran requires that the memory requirements of a program be determined entirely at compile time. This constraint, often viewed as a limitation by modern standards, was a strategic advantage for hardware access. In a static memory model, every variable, array, and constant is assigned a fixed, absolute memory address relative to a relocatable offset determined at load time. This eliminates the need for a runtime environment to manage a stack or a heap. When a Fortran program accesses a variable, the processor can use direct or absolute addressing modes, which are typically faster than the relative or indirect addressing required for stack-based variables.

The exclusion of recursion in the Fortran standard was a direct consequence of this static model. Since each subprogram's local variables occupied a fixed location in memory, calling a function recursively would overwrite the values from the previous invocation. Sacrificing recursion allowed compilers to perform aggressive global optimizations. The compiler knew exactly where every piece of data resided, enabling it to pre-load values into registers and keep them there throughout the execution of a loop, a task that becomes significantly more complex in a dynamic memory environment. This deterministic layout meant that the binary executable of a Fortran program was essentially a map of the machine's memory. Users could predict the exact cache behavior and memory bandwidth requirements of a program because there was no hidden memory management occurring in the background.

The common and equivalence statements provided a mechanism for low-level memory control that is rarely found in later high-level languages. These statements allowed programmers to explicitly define the physical overlap and alignment of data in memory. The common statement established named or unnamed blocks of memory that were shared across different subprograms. This was not merely a way to avoid passing arguments. It was a way to define the program's global memory map. By placing variables in a common block, the programmer ensured they were laid out in a contiguous sequence. This was critical for performance on architectures with limited memory, as it allowed for the creation of overlays, reusing the same physical memory for different parts of a program that did not need to be active simultaneously. 

The equivalence statement allowed two or more variables in the same subprogram to share the same memory location. While often used for memory conservation, its primary value for hardware access was in data type punning. By equivalencing a real variable with an integer array, a programmer could perform bitwise operations on the internal representation of a floating-point number. This was essential for implementing fast square root or trigonometric approximations by manipulating the exponent and mantissa bits directly, and compressing multiple smaller data types such as characters or 16-bit integers into a single 36-bit or 60-bit machine word. If a programmer combined common and equivalence in a way that violated memory alignment rules, the compiler would issue an error or insert padding bytes. This forced a level of architectural awareness that made Fortran code highly sympathetic to the processor’s memory bus and cache architecture.

The practice of mixing high-level Fortran code with assembly language was born out of necessity. During the mainframe era, programmers adopted a pragmatic, hybrid approach: they wrote the high-level structure and logic of a program in Fortran, but pushed performance-critical inner loops down into assembly. This allowed them to bypass the compiler when necessary and manually extract the maximum possible speed from the hardware. To make this integration seamless, systems programmers and compiler writers established strict subroutine linkage conventions, which served as the binary interfaces of their day. Writing assembly code that interfaced with Fortran required a deep understanding of the rules. The calling Fortran program constructed a parameter list in memory containing the contiguous addresses of the arguments. The address of this entire list was then passed to the assembly subroutine via a designated hardware register. The called assembly routine was strictly responsible for preserving the machine's state. The assembly routine had to save the caller's registers, establish its own base registers, and restore everything perfectly before returning control. If the assembly routine was a function, the result had to be left in a specific register.

The art of hand-coding assembly in symbiosis with Fortran ultimately faded. As hardware architectures evolved, such as the move from the IBM 7094 and CDC 6600 to the DEC-10, high-level Fortran code could be moved easily, but custom assembly routines required radical, expensive rewrites. This strategic liability drove the numerical computing community to concentrate low-level mathematical operations into standardized, highly optimized libraries that remain the bedrock of high-performance computing today. And crucially, by the Fortran 66 generation of compilers, they could often optimize the hardware's internal timing better than a human could, rendering hand-coded assembly largely unnecessary. Compiled Fortran becoming as efficient as custom assembly code was a milestone in the history of automation replacing manual human effort. Less human expertise and effort were needed, and software projects took less time, as Fortran compilers evolved.

Both assembly and Fortran had a concrete concept of floating-point numbers built into their foundations, but until Fortran 77 neither had a similar concept of text strings of characters. One of the difficult and fragile aspects of the assembly and Fortran combination was text manipulation. Before Fortran 77 introduced the formal character data type, handling text was a troublesome, hardware-dependent exercise that treated letters and symbols as simply groups of bits. Because early mainframe architectures like the CDC 6600 had no instructions for addressing individual characters, programmers had to write complex code to manually shift and mask bits just to isolate or extract a single character. Extracting the 7th through 14th characters of a 20-character string could require ten distinct machine instructions. Comparing two strings to see which came first alphabetically was considered a formidable task that could require fifty to sixty complex assembly instructions just to scan and evaluate the packed bits. This approach was fraught with difficulties and technical perils. In Fortran, the accepted convention was to handle text as Hollerith constants, named after the creator of the IBM punch card. Text data was in some sense still in the punch card era, though the physical punch cards themselves were less evident. Assembly routines were often written with hard-coded assumptions about these virtual punch cards, and the convention of six 6-bit characters packed into a single 36-bit word was still prevalent, meaning upper-case characters only and six character names. All of this directly impacted the design choices in the Decwar code. 

Decwar was a kind of culmination for the mainframe era in Texas and of simulation and modeling on big iron in Houston and Austin. And it is still alive and with us today, having survived some very surprising adventures and near-death experiences. Anyone today can easily boot up Tops-10, type \textit{r gam:decwar}, and be in the game. Even more startling is the fact that they can also freely work on the code itself using the original Teco, Macro-10, and Fortran-10, exactly as in 1978. All of this is purely old bits, installed from copies of the original DEC distribution tapes, except in the case of the game itself, which is a reconstruction of the UT Austin Source Distribution Tape based on the original manifest and code comments.

The game brought together many of the classic mainframe elements for one last hurrah before the personal computer era truly arrived in the eighties. There was no possibility of the Decwar code itself running on an eighties personal computer. For a microcomputer user, their machine could act as a video terminal, usually emulating a VT52 or VT100. Personal computers often in fact contained the exact same central processing chips used in commercial video terminals, which were microcomputers of a sort but not targeted at the personal computer market.

One of the classic mainframe elements genetically encoded within the game is that from the perspective of Tops-10, each player is simply a job, the same terminology and concept stretching back to batch jobs and punch card decks. For eighteen players, there are eighteen jobs, interleaved with each other by timesharing but directly descended from the old single-user batch jobs before timesharing. When none of the eighteen jobs is performing an action, nothing happens within the game. It is exactly like eighteen players sitting around a tabletop boardgame, or conducting a tabletop gaming session via physical mail in the old play by mail schemes. While no player is performing an action on the board, nothing happens in the game. It is turn-based at its core, with a turn occurring whenever a player makes an action on the game board. In fact, the term game board appears in the game’s source code. The game board is exactly the shared high segment chunk of magnetic-core memory statically owned by the game. The eighteen player jobs are capable of acquiring a lock on the shared memory and making a change to its contents. That is all. It is a direct extension of the two-player game with players taking turns at a teletype, but here eighteen players take turns accessing a shared memory.

This led directly to the role of the Romulan, to liven things up when there were only a few or just one player. It was extremely effective at raising the tension level, and also a vivid demonstration of the underlying reality of how the game worked. If you were the only player and were under attack by the Romulan, you could escape the stress by simply walking away from the keyboard. Nothing would happen as long as you didn’t type any commands. The Romulan had no dedicated job of its own. Instead, its actions and history emerged cooperatively from within the player’s jobs. In a sense, the player jobs collaboratively generated their own nemesis in the form of the Romulan. Whenever a player entered a command, their job could run the Romulan's subroutines and update the Romulan’s coordinates, energy, and shields, not to mention causing the Romulan to attack or transmit threats over the radio.

This adversary evolved from the 1974 single-player CDC 6600 Fortran version of the game, where developer Dave Matuszek created a computer-controlled opponent named The Thing. When the codebase was migrated to the DEC-10, the adversary was rebranded as the Romulan, but it kept the original question mark symbol used to denote The Space Thing. If the number of player jobs was below a threshold and there was not an active Romulan, each player job was capable of spawning a new one. Placing an automated agent into a real-time, shared-memory environment where players were constantly updating the global state led to some unintended technical challenges. Decwar players routinely stacked multiple abbreviated commands on a single input line, separated by slashes, for example \textit{MO RE 5 4/TO CO 3 R/PH CO EN}. In early versions, using the abbreviation R to target the Romulan caused internal parser collisions with a Klingon ship named the Raven. To handle this conflict, the Klingon ship was renamed the Manta. And because the Romulan's pathing algorithm initially computed simple straight-line vectors toward the nearest player ship, the cooperative updates would frequently steer it directly into black holes or within range of starbases. Developers had to patch the Romulan's logic with look-ahead checks to prevent it from navigating into danger. Even with these checks, human players could still cleverly force the Romulan into lethal hazards by using photon torpedoes and star nova explosions.

What really made the Romulan special was that it spoke, projecting an aura of sentience that fundamentally altered the tone of the game. For players in the early eighties, this was a psychological disruption, and they often found the Romulan’s ability to initiate direct, contextual communication thrilling and disturbing. It was more minimalistic and primitive and at the same time more primeval and powerful than the ritualized in-game encounters of today. It’s easy to identify people that actually played the game back in that era. They soon begin talking about how they’ll never forget their first encounters with the Romulan, and how they were instinctively convinced that the Romulan was a real and malevolent person at a remote terminal.

This communication was handled via radio messages, which, along with attack messages, formed dual dataflows within the code. These two dataflows were the dynamic core of both the shared high segment memory and of the overall game, and crucial to understanding the complete system. Both message types were stored in pre-allocated, static regions of the shared magnetic core memory. The radio message region immediately followed the attack message region, leading to a potential bug where running over the attack message region boundary corrupts the radio message data, and of course this has inadvertently been rediscovered and caused trouble during modern reconstruction efforts. The crucial point is that proper and sufficiently sized regions of core memory have to be allocated, otherwise the Decwar player jobs can unwittingly store attack or radio message data in the wrong region.

Focusing just on the radio message region, the layout of the messages form a linked list. The same is true for the attack message region and the following discussion applies to both equally. The data of each message contains links to the preceding and following messages. More importantly, each message contains a source identifier and destination identifiers, much like an email. Periodically, each player’s job checks if it is the destination for any messages. If so, it handles a radio message by displaying it on the player’s terminal. Then comes the critical phase, requiring a lock on the memory. The player-job must modify the message by removing itself as a destination, and if the message has no more destinations, remove the message from the linked list. Only one player-job can be allowed to do this at a time, so a lock has to be acquired first, and released once complete. 

The UT coders had to use Macro-10 assembly to invoke specific Tops-10 unimplemented user operations or uuo’s. Uuo’s acted as traps or interrupts that suspended user-level execution and called the Tops-10 monitor. The monitor would then assert exclusive access over the core memory segment, queuing any other player-jobs attempting to write to the lists until the holding job issued an unlocking uuo. In other words, in order to use the shared memory linked lists for attack and radio messages, player-jobs have to call Tops-10 directly via assembly code. This is the root reason that assembly code became essential for Decwar, and code in pure Fortran was simply no longer feasible. It unlocked the full potential of the DEC-10, and at the same time forever bound the code to Tops-10. 

There was very little separation between the game and the operating system, and little hope of moving the code to even a sibling environment such as the DEC-20 and its Tops-20 operating system. Even though Tops-20 ran on the same 36-bit hardware, it was derived from BBN's Tenex operating system rather than Tops-10, and it used a completely different system-call architecture known as jump to system or jsys instead of uuo’s. While Tops-20 did include a compatibility emulation library designed to intercept and translate old Tops-10 uuo’s into jsys calls, this emulator had severe limitations. The emulator was completely incapable of translating the direct physical segment locks and atomic inter-job synchronization routines that Decwar used to protect its message queues. Because the emulator relied on virtual memory page-mapping structures rather than static core segment locks, the Decwar binary simply could not run on Tops-20 without a complete rewrite of its underlying assembly code.

The cancellation of the PDP-10 product line by DEC in 1983, and the demise of DEC itself in 1998, seemed to permanently strand the UT Austin Decwar code on an obsolete architectural island. In another of its surprising near-death experiences, however, the UT code has survived and flourished, in an infinitely tougher and more survivable form, through the physical and digital preservation efforts led by Obsolescence Guaranteed. The mid 2020s saw a spectacular revitalization of this ecosystem. Rather than treating historical computing as dead artifacts meant only for museums, Obsolescence Guaranteed, an informal group of computer history hobbyists, had focused on creating computer time capsules. By building affordable, fully functional hardware replicas of the classic DEC lineup, including the PiDP-1, PiDP-8, PiDP-10, and PiDP-11, they allow modern users to directly experience the tactile and interactive realities of the mainframe era.

For Decwar, the PiDP-10 replica is one form of resurrection. The PiDP-10 is a scaled-down, desktop-sized physical reproduction of the original PDP-10 KA10 front panel. It features an active array of 74 functional switches and 124 indicator lamps, driven by modern LEDs, that accurately reflect the machine's internal state. Inside this console beats a dual-hearted system: a modern Raspberry Pi that runs a physical Linux kernel, side by side with a cycle-accurate PDP-10 emulator based on Bob Supnik and Richard Cornwell's Simh engine. Physical switch toggles on the front panel trigger interrupts on the Raspberry Pi, altering register states in the PDP-10, while memory writes are translated in real-time to illuminate the indicator lamps. 

A separate and even more permanent and accessible resurrection is represented by Docker. The modern reconstruction of the UT Austin Decwar Source Distribution Tape is a digital time capsule neatly packaged within a docker container. Docker excels at creating reproducible, automated build pipelines. A docker image encapsulates the underlying Simh emulator, the Tops-10 operating system, and the SDT virtual tape into a single, isolated package. This guarantees that anyone can instantly start up the living environment as it existed on the HRC DEC-10 in 1982, without worrying about local hardware dependencies or host operating system characteristics. In fact, running in the cloud is exactly the same as on local hardware. 

The preservation efforts go far beyond running a single standalone mainframe. The Arpanet Reconstruction Project aims to revive the ancestral Internet using old bits wherever possible. For UT Austin this would mean containers for the HRC DEC-10, Painter Hall DEC-20, and Arpanet IMP, with various connections through telnet and ftp links. Docker, especially using docker compose, is explicitly designed to orchestrate multi-node, networked applications. Instead of a user manually launching and configuring separate environments for mainframes and IMPs, a containerized setup can automatically start up each historical machine in its own isolated container and seamlessly manage the networking between them. Placing the Simh PDP-10 engine, historical tape images, and disk images inside a container shields users from the friction of modern software dependencies. They don't need to compile code or configure environments, but can simply start containers and immediately telnet into an old bits environment.
